% !TEX root = ../paper.tex

\section{Conclusion}
\label{sec:conclusion}

This paper studies how informal regulations and the discretion of individual political actors can influence housing supply within a city.
I focus on Chicago's aldermanic privilege, the informal power of each of the city's 50 aldermen to unilaterally block or approve development in their ward, and show that it has measurable effects on development activity in the city.

Using permit processing times for projects where aldermen are most likely to matter, I construct aldermanic stringency scores that isolate the alderman-specific component of regulatory delay after removing variation attributable to ward characteristics.
Using 2015 ward redistricting as identifying variation, I find that reassignment toward greater stringency reduces the annual number of high-discretion permits recorded by the city by about 8\% in the five years following the remap.
A boundary comparison shows that new construction is less dense on the more stringent side of ward boundaries, with multifamily FAR about 15\% lower and dwelling units per acre about 19\% lower.
I also find that listed rents are about 1.7\% higher on the more-stringent side after adjusting for observed unit and amenity differences, with suggestive evidence that home sale prices are about 1.6\% higher.
The supply-side designs provide causal evidence that variation in informal discretionary authority, beyond the formal rules on the books, affects housing production.
My results suggest that reforming the land-use policy in Chicago could affect housing production, but that effective reform would require addressing the power of individual discretion in the process along with the formal rules that have been the subject of most prior research.
